\chapterhead{1}{What a Mixture Is}{Parts that stay themselves when they are mixed}\label{ch1}

Trail mix, muddy water, a bowl of cereal and a pocketful of coins are all \textbf{mixtures}. A mixture has two or more parts that are put together but are not changed into something new. Each part is still what it was. Rice in a bowl of rice and lentils is still rice.

That means a mixture can usually be taken apart again. The big question of this book is how, and how to check that nothing went missing on the way.

\begin{tryit}{Sort by hand}
\textbf{You need:} a kitchen scale in grams, a bowl, three cups, about 100~g of rice, 60~g of lentils and 20 steel paper clips.
\begin{enumerate}[leftmargin=1.6em]
\item Weigh the rice, the lentils and the clips separately. Write each weight.
\item Put all three into the bowl and mix. Weigh the whole mixture. Write it.
\item Pick the three kinds apart by hand into the three cups.
\item Weigh each cup. Add the three weights. Compare the total with the mixture.
\end{enumerate}
\par\medskip
\renewcommand{\arraystretch}{1.6}
\begin{tabular}{|l|c|c|c|c|c|}
\hline
 & \textbf{Rice} & \textbf{Lentils} & \textbf{Clips} & \textbf{Whole mixture} & \textbf{Parts added}\\ \hline
\textbf{Before mixing (g)} & & & & & \\ \hline
\textbf{After sorting (g)} & & & & & \\ \hline
\end{tabular}
\end{tryit}
\adult{Steel clips and small beans can be swallowed by little children and pets. Work at a table and put everything away afterwards. Do not taste anything.}

\begin{worked}{Practice numbers, not from a real test}
Whole mixture on the scale: 180~g.\\
After sorting: rice 99~g, lentils 60~g, clips 20~g. Total $99+60+20=179$~g.\\
The parts are 1~g short of the mixture. A grain of rice rolled off the table, and the scale itself can be 1~g off. We call 1~g \emph{within the scale's error}.
\end{worked}

\begin{center}
\begin{tikzpicture}
  \node[bigcir,minimum size=2.6cm] (a) at (0,0) {};
  \node[cir,fill=tomato!30,inner sep=1pt] at (-0.45,0.35) {99}; \node[cir,fill=leaf!30,inner sep=1pt] at (0.45,0.35) {60}; \node[cir,fill=sky!40,inner sep=1pt] at (0,-0.45) {20};
  \node[valcir] (b) at (5.2,0) {179 g};
  \draw[popA] (a) -- node[above,note] {add the parts} (b);
  \node[note,anchor=west,text width=3.2cm,align=left] at (6.7,0) {\textbf{kept:} the weight of every part\\\textbf{set aside:} how the mixture looked};
\end{tikzpicture}
\end{center}

\begin{multix}[hmm]\multixbot[hmm]\tcblower
After I sorted my mixture the parts weighed 1 gram less than the whole. So one gram of stuff stopped existing. I will write that down. Should I?
\end{multix}

\trust{We weighed the whole mixture and each part, then added the parts and compared with the whole.}
{The parts came within 1~g of the whole, which is the size of the scale's own error.}
{That nothing at all was lost. A gap smaller than the scale's error cannot be told apart from no gap.}

\begin{keepbox}
We replaced a bowl of mixed things with three weights. \textbf{We kept} \wline[3.4cm] \textbf{and we set aside} \wline[3.4cm].
\end{keepbox}

\clearpage

\begin{puzzle}{Add the parts}
After sorting, rice was 98~g, lentils 61~g and clips 20~g. Total \wline[1.6cm]~g. The whole mixture weighed 180~g. How far off? \wline[1.6cm]~g
\end{puzzle}

\begin{puzzle}{Mixture or not?}
Write M for mixture and N for not a mixture.\\
(a) trail mix \wline[1cm] \quad (b) a spoon of pure salt \wline[1cm] \quad (c) muddy water \wline[1cm] \quad (d) one steel paper clip \wline[1cm]
\end{puzzle}

\begin{puzzle}{Too big a gap}
A 150~g mixture was sorted into 70, 50 and 25~g. Total \wline[1.6cm]~g. The gap is \wline[1.6cm]~g. Is that within the scale's 1~g error? \wline[2cm]\\
What would you do first? \wline[4cm]
\end{puzzle}

\begin{puzzle}{Spot the flaw}
Sam weighed the whole mixture on a kitchen scale and the parts on a bathroom scale. The parts added to more. Why is this test not trustworthy? \wline[6cm]
\end{puzzle}

\begin{puzzle}{Still themselves}
In a fruit salad, the grapes are still grapes. Is that true of every mixture? Explain. \wline[6cm]
\end{puzzle}

\clearpage

\begin{puzzle}{Draw the circles}
The parts of a mixture were 40, 25 and 15~g. Fill in the empty circle with the total.
\begin{center}
\begin{tikzpicture}
  \node[cir] at (0,0) {40}; \node[cir] at (1.4,0) {25}; \node[cir] at (2.8,0) {15};
  \node[valcir,dashed] (b) at (6.3,0) {}; \draw[popA] (3.6,0) -- node[above,note] {add} (b);
\end{tikzpicture}
\end{center}
\end{puzzle}

\begin{puzzle}{Biggest part}
Rice 99~g, lentils 60~g, clips 20~g. Which part weighs the most? \wline[2cm] Is it more than half of the 179~g? \wline[2cm]
\end{puzzle}

\begin{puzzle}{Your own mixture}
Name a mixture at home and its parts.\\
\wline[9cm]
\end{puzzle}

\begin{puzzle}{Rounded weights}
The scale showed 179~g for the parts. It showed 180~g for the whole. The smallest step on the scale is 1~g. Is the difference bigger than the smallest step? \wline[2cm]
\end{puzzle}

\begin{puzzle}{Missing grain}
A student says one grain of rice vanished. Write a more careful sentence. \wline[6cm]
\end{puzzle}
\clearpage
\thispagestyle{fancy}
\subhead{Mixture log}
Look for mixtures at home. Name the parts and say how you could take them apart.\par\medskip
\renewcommand{\arraystretch}{3.0}
\begin{tabularx}{\linewidth}{|>{\sffamily\small\bfseries}p{2.6cm}|X|X|}
\hline
Mixture & Parts & How I could separate them \\ \hline
 & & \\ \hline
 & & \\ \hline
 & & \\ \hline
 & & \\ \hline
 & & \\ \hline
\end{tabularx}
\clearpage
